\begin{center}
{\Large \bfseries \textcolor{lyred}{第六章 定积分的应用}}
\end{center}

\vspace{6pt}
{\large \bfseries \textcolor{lyred}{第6.1 节 定积分的元素法(微元法)}}

一般地,求一个与[
]
,a b 有关的量U ,若它关于区间具有可加性,即 
[
]
[
]
[
]
,
,
,
a b
a c
c b
U
U
U
=
+
,并且[
]
,x x
x
\forall 
+\Delta 
对应
[
]
()
(
)
,x x
x
U
U
f x
x
o
x
+\Delta 
\Delta 
=
=
\Delta +
\Delta 
,即 
()
dU
f x dx
=
,则
()
b
a
U
f x dx
=\int 
,其中
()
dU
f x dx
=
称为U 的\textcolor{lyred}{元素}. 
\textcolor{lyred}{例}.设一细棒位于[
]
0,3 之间,线密度
()
x
x
\rho 
=+
,求质量. 
解. [
]
[
]
,
0,3
x x
dx
\forall 
+
\subset 
,
()
dm
x dx
\rho 
=
,
()
(
)
M
x dx
x
dx
\rho 
=
=
+
=
\int 
\int 
. 
\textcolor{lyred}{例}.设有半径为R 的球体,体密度为
2r
\rho =
,求它的质量,其中r 为 
(1)点到球心的距离;(2)点到对称轴的距离; 
(3)点到过球心的某个平面的距离. 
解.(1)
R
M
r
r dr
R
\pi 
\pi 
=
\cdots 
=
\int 
; 
(2)
R
M
r
r
R
r dr
R
\pi 
\pi 
=
\cdots 
\cdots 
-
=
\int 
; 
(3)
(
)
R
M
r
R
r
dr
R
\pi 
\pi 
=
\cdots 
-
=
\int 
. 
\vspace{6pt}
{\large \bfseries \textcolor{lyred}{第6.2 节 定积分在几何学上的应用}}

\vspace{4pt}
{\bfseries \textcolor{lyred}{一.平面图形的面积}}

{\bfseries \textcolor{lyred}{1.直角坐标下的计算}}

{\bfseries \textcolor{lyred}{(1)一般方程}}

\textcolor{lyred}{情形1}.设D : a
x
b
\le 
\le 
,
()
()
y
x
y
y
x
\le 
\le 
,则
()
()
b
a
A
y
x
y
x
dx
=
-




\int 
. 
\textcolor{lyred}{情形2}.设D :c
y
d
\le 
\le 
,
()
()
x
y
x
x
y
\le 
\le 
,则
()
()
d
c
A
x
y
x
y
dy
=
-




\int 
. 
\textcolor{lyred}{例}.求
y
x
=
和
2y
x
=
所围图形的面积. 
解.视为X 型区域,则
(
)
A
x
x
dx
x
x


=
-
=
-
=




\int 
; 
视为Y 型区域,则
(
)
A
y
y
dy
y
y


=
-
=
-
=




\int 
. 
\textcolor{lyred}{例}.求
y
x
=
, y
x
=
和
x =
所围图形的面积. 
解.视为X 型区域,则
A
x
dx
x


=
-
=




\int 
; 
视为Y 型区域,
(
)
A
dy
y dy
y


=
-
+
-
=






\int 
\int 
. 
\textcolor{lyred}{例}.求
y
x
=
,
y
x
=
和
y =所围图形的面积. 
解.视为Y 型区域,
(
)
A
y
y dy
=
-
=
\int 
; 
视为X 型区域,
A
x
x
dx
x
dx




=
-
+
-
=








\int 
\int 
. 
\textcolor{lyred}{例}.求
y
x
=
与
y
x
=
-
所围图形的面积. 
解.
2,8
y
x
x
y
x

=
\Rightarrow 
=

=
-

,视为Y 型区域,则
A
y
y
dy
-


=
+
-
=




\int 
; 
视为X 型区域,则
(
)
(
)
A
x
x
dx
x
x
dx




=
--
+
-
-
=




\int 
\int 
. 
{\bfseries \textcolor{lyred}{(2)参数方程}}

\textcolor{lyred}{例}.求椭圆
x
y
a
b
+
=所围图形的面积. 
解.
(
)
sin
cos
sin
a
A
ydx
b
t d a
t
ab
tdt
ab
\pi 
\pi 
\pi 
=
=
\cdots 
=-
=
\int 
\int 
\int 
. 
\textcolor{lyred}{例}.求摆线
(
)
(
)(
)
sin
1 cos
x
a t
t
t
y
a
t
\pi 
=
-

\le \le 

=
-

与x 轴所围图形的面积. 
解.
(
)(
)
(
)
1 cos
sin
1 cos
a
A
ydx
a
t d at
a
t
a
t
dt
a
\pi 
\pi 
\pi 
\pi 
=
=
-
-
=
-
=
\int 
\int 
\int 
. 
\textcolor{lyred}{例}.求星形线
(
)
cos
sin
x
a
t
t
y
a
t
\pi 
=
\le \le 

=

所围图形的面积. 
解.
(
)(
)
(
)
sin
cos
sin
sin
a
A
ydx
a
t d a
t
a
t
t dt
a
\pi 
\pi 
\pi 
=
=
=
-
=
\int 
\int 
\int 
. 
{\bfseries \textcolor{lyred}{2.极坐标下的计算}}

\textcolor{lyred}{定理}.设\textcolor{lyred}{曲边扇形}D :\alpha 
\theta 
\beta 
\le 
\le 
,
()
\rho 
\varphi \theta 
\le 
\le 
,则它的面积
()
A
d
\beta 
\alpha 
\varphi 
\theta 
\theta 
=\int 
. 
\textcolor{lyred}{例}.求阿基米德螺线
(
)
a
\rho 
\theta 
\theta 
\pi 
=
\le 
\le 
与极轴所围图形的面积. 
解.
(
)
A
a
d
a
a
\pi 
\pi 
\theta 
\theta 
\theta 
\pi 


=
=
=




\int 
. 
\textcolor{lyred}{例}.求心形线
(
)(
)
1 cos
a
\rho 
\theta 
\theta 
\pi 
=
+
\le 
\le 
所围图形的面积. 
解.
(
)
(
)
1 cos
2cos
cos
A
a
d
a
d
a
\pi 
\pi 
\theta 
\theta 
\theta 
\theta 
\theta 
\pi 
=
\cdots 
+
=
+
+
=




\int 
\int 
. 
\textcolor{lyred}{例}.求双纽线
cos2
a
\rho 
\theta 
=
所围图形的面积. 
解.
,
,
4 4
\pi \pi 
\pi 
\pi 
\theta 




\in -
\cup 








,故
cos2
A
a
d
a
\pi 
\theta \theta 
=
\cdots 
=
\int 
. 
\textcolor{lyred}{例}.求圆
2 sin
\rho 
\theta 
=
和双纽线
cos2
\rho 
\theta 
=
所围图形的面积. 
解.
2 sin
sin
cos2
\pi 
\rho 
\theta 
\theta 
\theta 
\rho 
\theta 

=

\Rightarrow 
=
\Rightarrow 
=

=

,故
(
)
2 sin
cos2
A
d
d
\pi 
\pi 
\pi 
\theta 
\theta 
\theta \theta 
=
+
=
\int 
\int 
\pi 
\pi 




-
-
+
-
=
+












. 
\textcolor{lyred}{补充练习 }
1.求
sin
y
x
=
与
2x
y
\pi 
=
所围图形的面积. 
解.
sin
2 1
2 2
A
xdx
\pi 
\pi 
\pi 
\pi 






=
-
\cdots 
=
-
=
-










\int 
. 
2.求
y
x
=
与
y
x
=-
+
所围图形的面积. 
解.
y
x
x
y
x

=
\Rightarrow 
=\pm 

=-
+

,故
(
)
A
x
x
dx
-


=
-
+
-
=


\int 
. 
3.求
y
x
=
与
x
y
+
=
所围图形的面积. 
解.
y
x
y
x
y

=
\Rightarrow 
=-
+
=

,2 ,故
y
A
y
dy
-


=
-
-
=




\int 
. 
4.求曲线
ln
y
x
=
与两直线
y
e
x
=
+-
及
y =
围成平面图形的面积. 
解.
(
)
ln
e
e
e
A
xdx
e
x dx
+
=
+
+-
=
\int 
\int 
;或者,
(
)
y
A
e
y
e
dy
=
+-
-
=
\int 
. 
5.求
(
)
2 1
y
x
=
+
与
(
)
2 1
y
x
=
-
所围图形的面积. 
解.
y
y
y
A
dy
dy
y
-
-
-






=
-
-
-
=
-
=












\int 
\int 
. 
6.求心形线
(
)
1 cos
a
\rho 
\theta 
=
+
与圆
2 a
\rho =
所围图形的面积. 
解.
(
)
1 cos
a
a
\rho 
\theta 
\pi 
\theta 
\rho 
=
+

\Rightarrow 
=

=

,
(
)
1 cos
A
a
d
a
d
\pi 
\pi 
\pi 
\theta 
\theta 
\theta 






=
+
+
=










\int 
\int 
9 3
9 3
2 8
a
a
a
a
\pi 
\pi 
\pi 




+
-
=
-












. 
\vspace{4pt}
{\bfseries \textcolor{lyred}{二.体积}}

\textcolor{lyred}{定理(截面法)}.设立体\Omega 沿x 轴分布在a
x
b
\le 
\le 
的范围内,若与x 轴垂直的平面截 
该立体所得的截面面积为
()
A x ,则
()
b
a
V
A x dx
=\int 
. 
\textcolor{lyred}{推论}.
()
\{
\}
,0
D
a
x
b
y
f x
=
\le 
\le 
\le 
\le 
绕x 轴旋转所得旋转体的体积
()
b
x
a
V
f
x dx
\pi 
=\int 
. 
\textcolor{lyred}{注(薄壁桶法)}.设
a \ge 
,则上述图形绕y 轴旋转所得立体体积
()
b
y
a
V
xf x dx
\pi 
=
\int 
. 
\textcolor{lyred}{例}.求底面半径为r ,高为h 的圆锥体体积. 
解.
h
x
r
V
x
dx
r h
h
\pi 
\pi 


=
=




\int 
. 
\textcolor{lyred}{例}.求椭圆
x
y
b
a +
\le 分别绕x , y 轴旋转所得旋转体的体积. 
解.
a
x
a
x
V
b
dx
ab
a
\pi 
\pi 
-


=
-
=




\int 
;
b
y
b
y
V
a
dy
a b
b
\pi 
\pi 
-


=
-
=




\int 
. 
\textcolor{lyred}{例}.求圆
(
)
x
y
+
-
\le 绕x 轴旋转所得旋转体的体积. 
解.
(
)
(
)
x
V
x
dx
x
dx
x dx
\pi 
\pi 
\pi 
\pi 
-
-
-
=
+
-
-
-
-
=
-
=
\int 
\int 
\int 
. 
\textcolor{lyred}{例}.求
sin
y
x
=
与
2x
y
\pi 
=
所围图形绕x 轴旋转所得旋转体的体积. 
解.
sin
x
x
V
xdx
dx
\pi 
\pi 
\pi 
\pi 
\pi 
\pi 


=
\cdots 
-\cdots 
=




\int 
\int 
. 
\textcolor{lyred}{例}.求
y
x
=
与
x =
所围图形绕y 轴旋转所得旋转体的体积. 
解.
y
y
V
dy
\pi 
\pi 
\pi 
-




=
\cdots -
=








\int 
;或者,
1 2
2 2
y
V
x
xdx
\pi 
\pi 
=
\cdots 
=
\int 
. 
\textcolor{lyred}{例}.求
(
)
sin
y
x
x
\pi 
=
\le 
\le 
绕
x
\pi 
=
处的切线旋转所得旋转体的体积. 
解.
(
)
1 sin
V
x
dx
\pi 
\pi 
\pi 
\pi 
=
-
=
-
\int 
. 
\textcolor{lyred}{例}.求
y
x
=
与
x =
所围图形绕
y =-旋转所得旋转体的体积. 
解.
()
(
)
(
)
1 2
1 2
1 2
V
A x dx
x
x
dx
xdx
\pi 
\pi 
\pi 
\pi 
-


=
=
+
-
-
=
=




\int 
\int 
\int 
. 
\textcolor{lyred}{例}.求圆
x
y
x
+
\le 
绕
x =
旋转所得旋转体的体积. 
解.
(
)
(
)
V
y
y
dy
y dy
\pi 
\pi 
\pi 
-
-




=
-
-
-
-
-
+
-
=
-
=








\int 
\int 
. 
\textcolor{lyred}{例}.求摆线
(
)
(
)(
)
sin
1 cos
x
a t
t
t
y
a
t
\pi 
=
-

\le \le 

=
-

与x 轴所围图形分别绕x 轴, y 轴旋转所得 
旋转体的体积. 
解.
(
)
(
)
1 cos
1 cos
a
x
V
y dx
a
t
a
t dt
a
\pi 
\pi 
\pi 
\pi 
\pi 
=
=
-
\cdots 
-
=
\int 
\int 
; 
(
)
(
)
(
)
sin
1 cos
sin
a
y
V
xydx
a t
t
a
t d a t
t
a
\pi 
\pi 
\pi 
\pi 
\pi 
=
=
-
\cdots 
-
-
=




\int 
\int 
. 
\textcolor{lyred}{例}.求经过半径为R 的圆柱底面中心,与底面夹角为\alpha 的平面截圆柱所得体积. 
解.
(
)(
)
tan
tan
R
R
V
R
x
R
x
dx
R
\alpha 
\alpha 
-
=
-
-
=
\int 
. 
\textcolor{lyred}{例}.求两个半径为R 的直交圆柱体公共部分的体积. 
解.
()
(
)
R
R
V
A x dx
R
x
dx
R
=
=
-
=
\int 
\int 
. 
\textcolor{lyred}{补充练习 }
1.从半径为R 的球体上截取一个高为h 的球冠,求球冠的体积. 
解.
(
)
R
x
R h
h
V
R
x
dx
R
h
\pi 
\pi 
-


=
-
=
-




\int 
. 
2.求
y
x
=
与
y
x
=-
+
所围图形绕x 轴旋转所得旋转体的体积. 
解.
(
)
(
)
V
x
x
dx
\pi 
\pi 
-


=
-
+
-
=




\int 
. 
3.求
y
x
=
,
x =及x 轴所围图形绕y 轴旋转所得旋转体的体积. 
解.
y
V
y dy
\pi 
\pi 
\pi 
=
-
=
\int 
; 或者,
y
V
x x dx
\pi 
\pi 
=
\cdots 
=
\int 
. 
4.求圆(
)
x
y
-
+
\le 
绕y 轴旋转所得旋转体的体积. 
解.
(
)
(
)
y
V
y
dy
y
dy
y dy
\pi 
\pi 
\pi 
\pi 
-
-
-
=
+
-
-
-
-
=
-
=
\int 
\int 
\int 
. 
5.求
(
)(
)(
)
2 1
y
x
x
x
=
-
-
\le 
\le 
与x 轴所围图形分别绕x , y 轴旋转所得体积. 
解.
(
)(
)
x
V
x
x
dx
\pi 
\pi 
=
-
-
=
\int 
,
(
)(
)
y
V
x x
x dx
\pi 
\pi 
=
-
-
=
\int 
. 
\vspace{4pt}
{\bfseries \textcolor{lyred}{三.平面曲线的弧长}}

\textcolor{lyred}{定理}.设
()
()
: x
x t
L
y
y t
=

=

,
t
\alpha 
\beta 
\le \le 
,光滑,则
()
()
s
x t
y t
dt
\beta 
\alpha 
'
'
=
+
\int 
. 
\textcolor{lyred}{推论}.设光滑弧段
()(
)
:
L y
y x
a
x
b
=
\le 
\le 
,则
()
b
a
s
y
x dx
'
=
+
\int 
. 
\textcolor{lyred}{推论}.设极坐标下
()(
)
:
L \rho 
\rho \theta 
\alpha 
\theta 
\beta 
=
\le 
\le 
,则
()
()
s
d
\beta 
\alpha 
\rho 
\theta 
\rho 
\theta 
\theta 
'
=
+
\int 
. 
\textcolor{lyred}{例}.求摆线
(
)
(
)(
)
sin
1 cos
x
a
y
a
\theta 
\theta 
\theta 
\pi 
\theta 
=
-

\le 
\le 

=
-

的长度. 
解.
(
)
(
)
1 cos
sin
2 1 cos
2 sin 2
ds
a
a
d
a
d
a
d
\theta 
\theta 
\theta \theta 
\theta 
\theta 
\theta 
=
-
+
=
-
=
,故 
2 sin
2cos
s
a
d
a
a
\pi 
\pi 
\theta 
\theta 
\theta 


=
=
-
=




\int 
. 
\textcolor{lyred}{例}.设半圆弧线材
y
R
x
=
-
的线密度
k
y
\rho =
-
,求质量. 
解.对于(
)
,x y 处的一小段圆弧,
(
)
dy
dm
ds
k
y
dx
dx
\rho 


=
=
-
+
=




(
)
R
k
R
x
dx
R
x
-
-
-
,故
R
R
k
R
x
M
R
dx
kR
R
R
x
\pi 
-
-
-
=
=
-
-
\int 
. 
\textcolor{lyred}{例}.求阿基米德螺线
a
\rho 
\theta 
=
对应\theta 从0 到2\pi 一段的长度. 
解.
()
()
ds
d
a
a d
a
d
\rho 
\theta 
\rho 
\theta 
\theta 
\theta 
\theta 
\theta 
\theta 
'
=
+
=
+
=
+
,故 
(
)
ln 2
a
s
a
d
\pi 
\theta 
\theta 
\pi 
\pi 
\pi 
\pi 


=
+
=
+
+
+
+




\int 
. 
\vspace{6pt}
{\large \bfseries \textcolor{lyred}{第6.3 节 定积分在物理学上的应用}}

\vspace{4pt}
{\bfseries \textcolor{lyred}{一.作功问题}}

\textcolor{lyred}{例}.求在位于原点的电荷量为q 的点电荷所产生的静电场中将单位正电荷沿x 轴 
正方向从a 处移到b 处时电场力所作的功. 
解.由库伦定律,
()
q
F x
k x
\cdots 
=
,故
b
a
q
W
k
dx
kq
x
a
b


=
=
-




\int 
. 
\textcolor{lyred}{例}.把弹簧由原长拉长6cm 至少要作多少功? 
解.设平衡位置为原点,则
()
0.06
0.0036
0.0018
W
kxdx
k
k J
=
=
\cdots 
=
\int 
. 
\textcolor{lyred}{例}.圆柱形容器中有一定量气体,在等温条件下气体膨胀,把活塞从a 处推到b 处 
(
)
a
b
>
,求移动过程中气体压力所作的功. 
解.设容器横向摆放在[
]
0,a 上,底部在原点,由于PV
k
=
,故x 处作用在活塞上的 
压力
()
k
k
k
F x
PS
S
S
V
xS
x
=
=
=
=
,于是
()
ln
b
b
a
a
k
b
W
F x dx
dx
k
x
a
=
=
=
\int 
\int 
. 
\textcolor{lyred}{例}.设直径20cm ,高80cm 的圆柱体内充满压强为
cm
N
的气体,等温条件下 
把气体体积压缩到一半,问至少需要作多少功. 
解.设容器横向摆放在[
]
0,0.8 上,底部在原点,由于PV
k
=
, x 处作用在活塞上的 
压力
k
k
F
PS
S
V
x
=
=
=
,故
0.8
0.4
k
W
dx
x
=
=
\int 
()
ln 2
k
J ,由
P =
,
0.1
0.8
V
\pi 
=
\times 
\times 
, 
得
k
\pi 
=
,故
()
800 ln 2
W
J
\pi 
=
. 
\textcolor{lyred}{例}.将水从高5m ,底半径3m 的圆柱形桶中吸出,问至少需作多少功. 
解.以桶口为原点,垂直向下为正向建立x 轴,则将位于[
]
,x x
dx
+
之间的一薄层水 
吸出需作的功为
dW
x
gdV
x
g
dx
\rho 
\rho 
\pi 
=
\cdots 
=
\cdots 
\cdots 
,因此 
()
()
W
x
g
dx
g
J
J
\rho 
\pi 
\rho \pi 
\gamma \pi 
=
\cdots 
\cdots 
=
=
\int 
,其中
g
\gamma 
\rho 
=
为水的比重. 
\textcolor{lyred}{例}.将水从半径R 的半球形容器中吸出至少要作多少功? 
解.
()
(
)
A x
R
x
\pi 
=
-
,故
(
)
R
W
x
g
R
x
dx
g R
\rho 
\pi 
\rho \pi 
=
\cdots 
\cdots 
-
=
\int 
. 
\textcolor{lyred}{例}.设半径为R 的匀质球浮在水中,球的上部与水面相切,若将球从水中取出,问 
至少需作多少功. 
解.
(
)
(
)
(
)(
)
R
R
W
R
x
g
R
R
x
dx
g
R
x
Rx
x
dx
\rho \pi 
\rho \pi 


=
-
\cdots 
-
-
=
-
-
=


\int 
\int 
(
)
(
)
R
R
g
R
Rx
x
dx
g
x
Rx
x
dx
gR
\pi 
\rho \pi 
\rho \pi 
\rho 
-
-
-
=
\int 
\int 
,注意\rho 
\rho 
=
球
水. 
\vspace{4pt}
{\bfseries \textcolor{lyred}{二.水压力}}

\textcolor{lyred}{例}.横放的圆柱形水桶内盛有半桶水,桶的底半径为R ,求桶的一端所受的压力. 
解.以水面为原点,垂直向下为正向建立x 轴,则对位于[
]
,x x
dx
+
之间的窄条, 
dF
gx
R
x dx
\rho 
=
\cdots 
-
,故
R
F
g x R
x dx
gR
\rho 
\rho 
=
-
=
\int 
. 
\textcolor{lyred}{例}.一底边长8m ,高6m 的等腰三角形铁板,垂直地浸没在水中,顶在上底在下, 
顶离水面3m ,求它所受压力. 
解.以水面为原点,垂直向下为正向建立x 轴,则由三角形的相似, 
()
()
(
)
l x
x
l x
x
-
=
\Rightarrow 
=
-
,故
(
)
F
gx
x
dx
g
\rho 
\rho 
=
\cdots 
-
=
\int 
(牛). 
\textcolor{lyred}{例}.设横放的椭圆柱型水箱高1.5米,宽2 米,求当水箱装满水时,它的一个端面所 
受到的水压力. 
解.设椭圆的方程为
0.75
y
x +
=,故
(
)
0.75
0.75
0.75
F
g
y
xdy
\rho 
-
=
-
\cdots 
=
\int 
(
)
0.75
0.75
0.75
1.5
0.75
0.75
y
g
y
dy
g
g
\rho 
\rho 
\pi 
\rho \pi 
-
-
-
=
\cdots 
\cdots 
=
\int 
(牛). 
\vspace{4pt}
{\bfseries \textcolor{lyred}{三.引力}}

\textcolor{lyred}{例}.设一单位质点位于长2m,线密度为
kg/m
\mu 
的均匀细杆的中点上方1m 处,求 
细杆对该质点的引力. 
解.
cos
y
dx
dF
dF
k
x
x
\mu 
\alpha 
\cdots 
=
\cdots 
=
\cdots 
+
+
,故
(
)
3 2
y
k dx
F
k
x
\mu 
\mu 
-
=-
=-
+
\int 
(牛), 
由对称性,
x
F =
(牛). 
\textcolor{lyred}{例}.设圆弧型材的半径R ,中心角2\varphi ,线密度\mu ,求它对圆心处单位质点的引力. 
解.圆心为极点,对称轴为极轴建立极坐标系,由对称性,
y
F =
, 
cos
cos
x
Rd
k
dF
k
d
R
R
\mu 
\theta 
\mu 
\theta 
\theta \theta 
\cdots 
=
\cdots 
=
,故
cos
sin
x
k
k
F
d
R
R
\varphi 
\varphi 
\mu 
\mu 
\theta \theta 
\varphi 
-
=
=
\int 
. 
\textcolor{lyred}{例}.设星形线
cos
sin
x
a
t
t
y
a
t
\pi 
=


\le \le 



=



上每一点处的线密度等于该点到原点距离的 
立方,求星形线对原点处单位质点的引力. 
解.
cos
cos
sin
x
x
r ds
dF
dF
k
kxds
ka
t
tdt
r
r
\alpha 
\cdots 
=
\cdots 
=
\cdots 
=
=
,故 
cos
sin
x
F
ka
t
tdt
ka
\pi 
=
=
\int 
,由对称性,
y
F
ka
=
. 
\textcolor{lyred}{补充练习 }
1.设装满水的容器内壁由曲线
(
)
3 0
y
x
x
=
\le 
\le 
(单位:米)绕y 轴旋转而成,将水 
吸到容器上方2 米处至少需要作多少功? 
解.
(
)
()
W
y
g
x dy
g
y
y
dy
g
J
\rho 
\pi 
\rho \pi 
\rho \pi 


=
-
\cdots 
\cdots 
=
-
=




\int 
\int 
. 
2.设一部挖掘机从30m 深的井中挖出污泥,设抓斗自身重40kg ,每次抓起的污泥 
重200kg ,且在提升过程中污泥以2kg/s 的速度从抓斗缝隙中漏掉,提升的速度为 
3m/s ,问抓斗每次抓满污泥后提升至井口至少需作多少功? 
解.设提升过程为时间段[
]
0,10 ,则在[
]
,t t
dt
+
间,克服重力作功为 
(
)
dW
t
g
dt
=
+
-
\cdots 




,故
(
)
()
W
t
g
dt
g J
=
+
-
\cdots 
=




\int 
. 
\textcolor{lyred}{注}.若考虑到5kg/m的缆绳,则
(
)
(
)
5 30
dW
t
t
g
dt
=
+
-
+
-
\cdots 




,于是 
(
)
(
)
()
5 30
W
t
t
g
dt
g J
=
+
-
+
-
\cdots 
=




\int 
. 